% Standalone note. From the repository root, compile with:
% latexmk -pdf -outdir=.vscode/latex-build/notes notes/stackelberg-counterexample.tex
\documentclass[11pt]{article}
\usepackage[T1]{fontenc}
\usepackage{times}
\usepackage{amsmath,amssymb,amsthm}
\usepackage[margin=1in]{geometry}
\usepackage{xspace}

\newcommand{\X}{\mathcal X}
\newcommand{\Y}{\mathcal Y}
\newcommand{\C}{\mathcal C}
\renewcommand{\H}{\mathcal H}
\newcommand{\V}{\mathcal V}
\newcommand{\Lo}{\mathcal L}
\renewcommand{\Re}{\mathcal R}
\newcommand{\Mi}{\mathcal M}
\newcommand{\hy}{\hat y}
\newcommand{\vy}{\tilde y}
\newcommand{\ind}{\mathbf 1}
\newcommand{\dt}{\texttt{Proposer}\xspace}
\newcommand{\df}{\texttt{Responder}\xspace}

\title{A regret-based Stackelberg equilibrium with linear prediction mistakes}
\author{}
\date{}

\begin{document}
\maketitle

We give a verifier class where cooperation makes at most one prediction
mistake, but a Stackelberg equilibrium for the debaters' individual regrets
has linearly many expected mistakes. The \df can achieve zero regret
without learning the true label. Its responses then leave the \dt dependent
on guessing a rare certificate. A different, cooperative query reveals the
true label immediately.

We use a Bayesian Stackelberg game. Both players know a common prior over
verifiers, and their objectives are expected individual regrets under this
prior. The verifier is sampled once and remains fixed. In particular, the
equilibrium claim below concerns this common-prior game, rather than
separately taking a worst case over verifiers for each player.

\section{Verifier class and feedback}

Let $\X=\{1,2,\ldots\}$, $\Y=\{0,1\}$, and
$\C=\{0,1,2,\ldots\}$. For each $x\in\X$, define the finite set
\[
S_x=\{3,4,\ldots,4x(x+1)+2\},
\qquad |S_x|=4x(x+1).
\]
For $b\in\{0,1\}$ and a function $s$ satisfying $s(x)\in S_x$ for every
$x$, define
\[
v_{b,s}(x,p,r)=
\begin{cases}
b, & p=s(x)\text{ or }r=b,\\
1-b, & \text{otherwise}.
\end{cases}
\]
The verifier class and hypothesis class are
\[
\V=\{v_{b,s}:b\in\{0,1\},\ s(x)\in S_x\text{ for all }x\},
\qquad
\H=\{h_0,h_1\},\quad h_b(x)=b.
\]
Each $v_{b,s}$ validates exactly $h_b$. The certificate $p=s(x)$ forces
output $b$ against every response. Conversely, $r=b$ always forces output
$b$, so no proposer certificate can force $1-b$ against every response.
If $p\ne s(x)$, a response different from $b$ gives output $1-b$.
Thus $s(x)$ is the unique winning proposer certificate.

The prior samples $b$ uniformly from $\{0,1\}$ and, independently,
samples each $s(x)$ uniformly from $S_x$. Neither player sees the sampled
$b$ or $s$. Their private random coins are independent of these samples.

Fix a finite horizon $T$ and present the distinct inputs $x_t=t$. In each
round, the \dt announces $\hy_t$ and $p_t$. After observing both, the \df
chooses $r_t$. Both players then observe
$\vy_t=v_{b,s}(t,p_t,r_t)$. Their policies may depend on the complete
public history and their private randomness. They have no additional
queries to the verifier, and the true label is not separately revealed.

\section{Regrets and the Stackelberg objective}

Prediction mistakes compare the proposed label with the true label:
\[
\Mi(T)=\sum_{t=1}^T\ind[\hy_t\ne b].
\]
The individual contest losses are
\[
\Lo_T(T)=\sum_{t=1}^T\ind[\vy_t\ne\hy_t],
\qquad
\Lo_F(T)=\sum_{t=1}^T\ind[\vy_t=\hy_t].
\]
Every round is achievable for a \dt that knows the verifier and may choose
both its label and its certificate. Hence its regret benchmark is zero,
and $\Re_T(T)=\Lo_T(T)$.

A round is unachievable for the \df exactly when the proposed label is
correct and the proposer certificate is winning. Define
\[
U_F(T)=\sum_{t=1}^T\ind[\hy_t=b,\ p_t=s(t)],
\qquad
\Re_F(T)=\Lo_F(T)-U_F(T).
\]
Every unachievable round contributes one to $\Lo_F$. Consequently each
round's contribution to $\Re_F$ is nonnegative. We do not subtract a round
merely because $p_t=s(t)$ when the proposed label is wrong.

Let $A$ and $B$ denote the complete history-dependent policies of the \dt
and \df. The \dt commits to $A$ before play, and the \df observes this
policy before choosing $B$. The commitment does not reveal private random
coins or the sampled verifier. Write $\Re_i(A,B;T)$ for the resulting
regret. A follower best-response map $B^*$ and a leader policy $A^*$ form
a Stackelberg equilibrium if
\begin{align*}
B^*(A)&\in\operatorname*{arg\,min}_{B}
    \mathbb E[\Re_F(A,B;T)] \qquad\text{for every }A,\\
A^*&\in\operatorname*{arg\,min}_{A}
    \mathbb E[\Re_T(A,B^*(A);T)].
\end{align*}
All expectations include the initially sampled verifier and the players'
randomness. Both objectives are individual regrets, rather than raw
loss for one player and regret for the other.

\section{A response with zero regret}

Consider the responder policy $B_0$ that always chooses
\[
r_t=1-\hy_t.
\]
This policy achieves $\Re_F(T)=0$ for every proposer policy and every fixed
verifier. If $\hy_t\ne b$, then $r_t=b$ and the \df wins. If $\hy_t=b$
but $p_t\ne s(t)$, then $r_t\ne b$, the verifier outputs $1-b$, and the
\df again wins. The only remaining case has $\hy_t=b$ and $p_t=s(t)$.
The \df loses in this case, but the round is unachievable and the loss is
subtracted.

Since responder regret is nonnegative, $B_0$ is an exact best response to
every proposer commitment. Moreover, unless the \dt guesses the correct
certificate, the feedback is
\[
p_t\ne s(t)
\quad\Longrightarrow\quad
(r_t,\vy_t)=(1-\hy_t,1-\hy_t).
\]
This feedback has the same form for both possible values of $b$.

\section{Linear prediction mistakes}

Fix any proposer policy $A$. Couple two runs against $B_0$, one with $b=0$
and one with $b=1$, using the same function $s$ and the same proposer
random coins. Until the first round with $p_t=s(t)$, the two runs have
identical histories and therefore identical proposed labels and
certificates.

Let $\tau$ be this first certificate hit, with $\tau=\infty$ if there is
none. Conditional on no earlier hit, the fresh value $s(t)$ is still
uniform on $S_t$ and independent of the current certificate choice.
Therefore
\[
\Pr(\tau=t\mid\tau\ge t)\le\frac{1}{4t(t+1)}.
\]
Using the union bound and a telescoping sum gives
\[
\Pr(\tau\le T)
\le\sum_{t=1}^T\frac{1}{4t(t+1)}
=\frac14\left(1-\frac{1}{T+1}\right)
\le\frac14.
\]
Let $\Mi_0(T)$ and $\Mi_1(T)$ be the mistake counts in the two coupled
runs. On $\{\tau>T\}$, the same prediction is made in both worlds each
round, so $\Mi_0(T)+\Mi_1(T)=T$. Averaging over the fair bit $b$ yields
\[
\mathbb E[\Mi(T)]
=\frac12\mathbb E[\Mi_0(T)+\Mi_1(T)]
\ge\frac T2\Pr(\tau>T)
\ge\frac{3T}{8}.
\]
This holds for every proposer policy, including one that optimally
minimizes its own regret against $B_0$.

\section{Existence of an equilibrium and follower ties}

Set $B^*(A)=B_0$ for every leader commitment $A$. This is a best-response
map by the zero-regret argument. For each finite horizon, choose $A_T^*$
to minimize $\mathbb E[\Re_T(A,B_0;T)]$.

Such a minimizer exists. Against $B_0$, all proposer certificates outside
$S_t$ are equivalent to $p_t=2$, since the verifier only checks whether
$p_t=s(t)$ and $B_0$ does not depend on $p_t$. The relevant action and
observation sets are therefore finite at each of the finitely many
rounds. Backward induction gives an optimal policy. Thus
$(A_T^*,B^*)$ is a Stackelberg equilibrium satisfying
\[
\Re_F(T)=0,
\qquad
\mathbb E[\Mi(T)]\ge\frac{3T}{8}.
\]
The optimal leader policy may depend on $T$, but the class, prior, and
input sequence are the same for every horizon.

The mistake lower bound also survives other choices among exact follower
best responses. Every best response has expected regret zero, and so
cannot incur a positive regret increment with positive probability.
Inductively, before a certificate hit, both values of $b$ have positive posterior
probability. Also, the fresh certificate misses $s(t)$ with positive
probability. If the \df chose any $r_t\ne1-\hy_t$, then on the event
\[
b=1-\hy_t,\qquad p_t\ne s(t),
\]
the verifier would output $\hy_t$. This would be an avoidable responder
loss, contradicting zero expected regret. Hence every exact follower
best response must use $r_t=1-\hy_t$ before the first hit. The same
coupling proof applies to every such response. Favorable tie-breaking
therefore cannot remove the prediction-mistake lower bound.

\section{Cooperative learning}

Cooperatively, choose $p_1=r_1=2$. Since $s(1)\ge3$ and
$b\in\{0,1\}$, this query gives
\[
v_{b,s}(1,2,2)=1-b.
\]
Both players can recover $b$ from the output. Thereafter the \dt always
predicts $b$, so $\Mi(T)\le1$ for every fixed verifier. The \df can also
choose $r_t=b$ thereafter to make the verifier itself output the true
label, regardless of the proposer certificate.

Thus the same verifier class permits learning the label with one
cooperative query, while exact regret-based Stackelberg play can keep it
hidden. The expectation in the equilibrium lower bound is over the
common prior; it is not a claim that every fixed verifier is hard. For
each horizon and policy pair covered by the bound, averaging does imply
that some fixed verifier has at least $3T/8$ expected mistakes over the
players' randomness.

The use of exact best responses matters. The diagnostic query $p=r=2$
incurs positive expected responder regret, so it is excluded by the
zero-regret alternative. An additive constant tolerance would permit
such a query. The argument above does not rule out accurate approximate
equilibria.

\end{document}
